% fonts.tex — настройка шрифтов для XeLaTeX + русский язык.
% Предпочитаем PT-семейство, но явно фолбэчим на шрифты из TeX Live.

\usepackage{fontspec}
\defaultfontfeatures{Ligatures=TeX}

\newcommand{\paymentsbookmainfont}{PT Serif}
\IfFontExistsTF{PT Serif}{}{%
  \IfFontExistsTF{Libertinus Serif}{%
    \renewcommand{\paymentsbookmainfont}{Libertinus Serif}%
  }{%
    \renewcommand{\paymentsbookmainfont}{TeX Gyre Termes}%
  }%
}

\newcommand{\paymentsbooksansfont}{PT Sans}
\IfFontExistsTF{PT Sans}{}{%
  \IfFontExistsTF{Libertinus Sans}{%
    \renewcommand{\paymentsbooksansfont}{Libertinus Sans}%
  }{%
    \renewcommand{\paymentsbooksansfont}{TeX Gyre Heros}%
  }%
}

\newcommand{\paymentsbookmonofont}{PT Mono}
\IfFontExistsTF{PT Mono}{}{%
  \IfFontExistsTF{Liberation Mono}{%
    \renewcommand{\paymentsbookmonofont}{Liberation Mono}%
  }{%
    \renewcommand{\paymentsbookmonofont}{Latin Modern Mono}%
  }%
}

%% Основной шрифт (с засечками) — хорошо читается в печатном виде.
\setmainfont{\paymentsbookmainfont}[Numbers=OldStyle]

%% Шрифт без засечек — для подписей, колонтитулов и боковых заметок.
\setsansfont{\paymentsbooksansfont}

%% Моноширинный шрифт — для листингов кода, hex-дампов и идентификаторов.
\setmonofont{\paymentsbookmonofont}[Scale=0.88]

%% Явные объявления кириллических шрифтов для polyglossia.
\newfontfamily\cyrillicfont{\paymentsbookmainfont}[Script=Cyrillic]
\newfontfamily\cyrillicfontsf{\paymentsbooksansfont}[Script=Cyrillic]
\newfontfamily\cyrillicfonttt{\paymentsbookmonofont}[Scale=0.88]
